\section{Discussion and conclusion}
\label{sec:discussion}

This work develops a variational analysis of finite-pool Chebyshev greedy approximation for nonlinear PDE energies. Under the atomic source condition, directional smoothness, and uniform positive selection strength, the ideal continuous CGA has an \(O(m^{-1})\) energy bound. For finite-pool CGA, the one-step estimate quantifies the contributions of dictionary coverage, score evaluation, radial stationarity, and coefficient optimization. Relative accuracy conditions preserve the same energy order.

For the pure, regularized, and reaction \(p\)-energies with \(p\ge2\), the energy gap controls both the natural \(V\)-distance and the Sobolev error. An energy exponent \(\alpha\) yields the metric exponents \(\alpha/2\) and \(\alpha/p\), respectively. These relations translate greedy energy estimates into error measures adapted to the nonlinear PDE geometry.

Local Hessian transfer and finite-trajectory projection estimates connect the descent analysis with approximation theory. The Hessian argument relates nonlinear selection to weak OGA in a frozen Hilbert metric. The finite-trajectory argument uses approximate source representations and normalized innovations to bound errors in the generated nested spaces. For quartic energies, the exact expansion converts these projection estimates into nonlinear energy bounds.

The one-dimensional pure and regularized \(p=4\) experiments illustrate the finite-trajectory estimates. On \(N=8,\ldots,32\), the source and quartic comparison bounds cover the computed energy errors. The prescribed benchmark \(\alpha_{\rm bench}=6\) gives comparison powers of \(6\), \(3\), and \(3/2\) for energy, natural distance, and Sobolev error. With the calibration constants fixed, the normalized-innovation condition holds on 27 of 32 continuation transitions for the pure model and 28 of 32 for the regularized model.

At the five prescribed comparison endpoints, CGA has smaller Sobolev errors than the metric-valid RFM medians and smaller errors than the P3 FEM references in four cases. These comparisons quantify the approximation accuracy obtained by adaptive greedy selection at matched coefficient counts.

The finite-pool implementation requires a balance between directional coverage and the costs of candidate evaluation and repeated coefficient optimization over growing trial spaces. These computational demands motivate residual-guided pool enrichment, preconditioned coefficient solvers, and adaptive optimization tolerances informed by the inexact-descent estimates. On the theoretical side, further questions concern the uniform control of the local Hessian-transfer and quartic-comparison conditions as the approximation spaces grow, and the construction of error estimators that can be evaluated without knowledge of the exact solution. The present numerical study focuses on smooth manufactured solutions in one and two dimensions. Extending these experiments to higher-dimensional problems and solutions of limited regularity, together with systematic comparisons at matched computational budgets, would clarify the scalability of CGA and the relation between its approximation accuracy and computational cost.

Together, the descent estimates, natural-distance relations, and projection bounds provide a quantitative framework for adaptive ridge approximation of nonlinear PDE energies. The numerical results demonstrate how finite-pool greedy selection constructs accurate trial spaces for the representative linear, semilinear, and \(p\)-structure problems.
